\documentclass[10pt,a4paper]{amsart}
\usepackage[margin=19mm]{geometry}
\usepackage{amsmath,amssymb,mathtools}
\usepackage[scr=rsfs,cal=euler]{mathalfa}
\usepackage{xfrac}
\usepackage[unicode,hidelinks,bookmarks=false]{hyperref}
\newcommand{\ie}{i.e.\ }
\newcommand{\cf}{cf.\ }
\newcommand{\resp}{respectively\ }
\newcommand{\Subsection}{\subsection}
\providecommand{\nobreakdash}{\nobreak\hbox{-}}
\setlength{\emergencystretch}{2em}
\multlinegap0pt
\usepackage{amssymb}
\newtheorem{rem}{Remark}
\newtheorem{thm}{Theorem}
\newtheorem{prop}{Proposition}
\newtheorem{lem}{Lemma}
\newcommand{\CC}{\mathbb C}
\newcommand{\RR}{\mathbb R}
\title{Simply connectedness of spaces of tight frames}
\author{Augustin-Liviu Mare}
\date{Source: arXiv:2508.02484v2 (CC BY 4.0)}
\begin{document}
\maketitle

\noindent\emph{Source and licence.} Simply connectedness of spaces of tight frames. Version arXiv:2508.02484v2 (CC BY 4.0), distributed by arXiv under the Creative Commons Attribution 4.0 International licence, http://creativecommons.org/licenses/by/4.0/, as recorded in the arXiv licence metadata for this exact version and shown on the abstract page https://arxiv.org/abs/2508.02484v2. Full licence notice: \texttt{licenses/arxiv-2508.02484-CC-BY-4.0.txt}.\par\smallskip

\noindent\textit{Editorial scope.} Counted unit: the article's prop key (source0.tex lines 397--399). The article's complete original proof of it is delivered with it (source0.tex lines 401--425, one \texttt{proof} environment of 25 lines). No mathematics is written, repaired or re-derived: the statement and the proof are the article's own lines, in the article's own notation, and no retained line is substituted or re-flowed.  Retained uncounted as the source-local context the proof needs: lines 159--160; lines 312--317; lines 321--339; lines 332--332; lines 369--374; lines 389--391.  Nothing was omitted from any retained span, so the package carries no omission record.  The retained lines cite 1 works, whose entries are transcribed verbatim from the article's own bibliography source in the e-print and delivered with short editorial labels recorded in the item's reference mapping.  No retained line contains a figure environment, an image-inclusion command or drawing code, so no image is delivered and the package declares no asset dependency.  Editorial formatting only, in addition: an AMS \texttt{amsart} preamble that restates the article's own declarations and macros used by the retained lines, namely the environments \texttt{rem}, \texttt{thm}, \texttt{prop}, \texttt{lem} on separate counters, together with the standard packages those lines need.  The retained environments print in this document as Remark 1, Theorem 1, Proposition 1, Lemma 1 in order of appearance, whereas the article prints the same items with the numbering of its own full document.  The complete native source of the article as distributed in the arXiv e-print for this exact version is bundled unchanged as \texttt{sources/182674/source0.tex}.
\begin{equation} \label{hypothesis} d_{i_1} +\cdots + d_{i_{n-k}}\ge 1 \ {\it 
 for \ any \ pairwise  \  distinct } \ i_1, \ldots, i_{n-k} \in \{1, \ldots, n\}, \end{equation} then ${\mathcal F}_{n,k}^d$

\begin{rem} \label{unif} Assume $a=(x, \ldots, x)$, that is, all components of $a$ are equal to each other. Then 
$h_a(V) = \langle \pi_V, x{\rm I}_n\rangle = {\rm tr}(x\pi_V)=xk$, which means that the function $h_a$ is constant. 
Thus all $V\in {\rm Gr}_k(\CC^n)$ are critical points for $h_a$. This is also confirmed by the theory above, since 
we clearly have $\ell =1$ and for any $u\subset \{1, \ldots, n\}$ with $k$ elements, $M_u^a$ is the whole 
${\rm Gr}_k(\CC^n)$. 
\end{rem}

\begin{thm}\label{stratif} \begin{itemize} 
\item[(i)] The critical set of $f_{}$ is the union of all (non-empty) intersections of the form
$\mu^{-1}(d+a) \cap gM_{u}^{{a}^\sigma}$ where $a\in \RR^n$  and $u\subset \{1, \ldots, n\}$ with $k$ elements.
The intersection is non-empty  only for finitely many  $a$ and $u$ as above. For any such $a$ and $u$, there is a submanifold
$\Sigma_{a, u}$  of ${\rm Gr}_k(\CC^n)$ which contains $\mu^{-1}(d+a) \cap gM_{u}^{{a}^\sigma}$
as a deformation retract. These submanifolds induce a partition 
\begin{equation}\label{stratsigma}{\rm Gr}_k(\CC^n)= \bigsqcup \Sigma_{a, u}.\end{equation}
The codimension of $\Sigma_{a,u}$ in
${\rm Gr}_k(\CC^n)$ is equal to the index of $h_a$ along $gM_{u}^{{a}^\sigma}$, that is, the codimension of
$gS_{u}^{a^\sigma}$.
\item[(ii)] The partition (\ref{stratsigma}) is a stratification in the sense of \cite[Def.~2.11]{Ki}. More precisely, one can label the submanifolds as $\{\Sigma_\beta  \mid \beta\in B\}$, where $B$ is a finite set endowed with a strict partial order ``$ \ >$" such that for any $\beta \in B$ the topological closure of $\Sigma_\beta$ satisfies
\begin{equation}\label{osig} \overline{\Sigma}_\beta \subseteq \bigcup_{\gamma \ge \beta} \Sigma_\gamma.\end{equation} 
\item[(iii)] There is exactly one stratum $\Sigma_{a,u}$ of codimension zero, namely 
the one corresponding to $a=(0, \ldots, 0)$ 
and $u$  any subset with $k$ elements of\footnote{See also Rem.~\ref{unif}. Also note the general fact that if $a=(a_1, \ldots, a_n)$ such that $\mu^{-1}(d+a)$ is non-empty, then $a_1+\cdots + a_n=0$, because the sum of the components of both $d$ and $d+a$ is equal to $k$.}
$\{1, \ldots, n\}$.  Moreover,  $M_{u}^{(0, \ldots, 0)}$ is the whole ${\rm Gr}_k(\CC^n)$ and hence 
$\Sigma_{(0, \ldots, 0),u}$ contains $\mu^{-1}(d)$ as a deformation retract. In terms of the notations from point (ii), this stratum is of the form $\Sigma_{\beta_0}$, where $\beta_0\in B$ is the greatest element.
\end{itemize} 
\end{thm} 

\begin{equation} \overline{\Sigma}_\beta \subseteq \bigcup_{\gamma \ge \beta} \Sigma_\gamma.\end{equation} 

\begin{prop} \label{prop1ton}
Assume $d\in \Delta_{n,k}$ satisfies  assumption (\ref{hypothesis}). If $a\in \RR^n$ and $u\subset \{1, \ldots, n\}$ with $k$ elements are such that 
$\mu^{-1}(d+a) \cap g M_{u}^{{a}^\sigma}$ 
is non-empty and the codimension of $\Sigma_{a,u}$ is strictly greater than 0
then the codimension of $\Sigma_{a,u}$  is at least equal to 4.
\end{prop}

\begin{lem}  \label{cs} {\rm (B.~Williams)} Let $M$ be a connected manifold and $Z\subset M$ a closed submanifold of codimension $d\ge 2$. Then the inclusion map 
$M\setminus Z \hookrightarrow M$ is $(d-1)$-connected.  
\end{lem}

\begin{prop} \label{key} If $d$ satisfies assumption (\ref{hypothesis}), then the inclusion map
$\mu^{-1}(d) \hookrightarrow  {\rm Gr}_k(\CC^n)$ is 3-connected.
\end{prop} 

\begin{proof}
The map mentioned in the proposition can be decomposed as 
$$\mu^{-1}(d) \hookrightarrow \Sigma_{(0,\ldots, 0),u} \hookrightarrow {\rm Gr}_k(\CC^n),$$
where $u\subset \{1, \ldots n\}$ is an arbitrary subset with $k$ elements. 
By Thm.~\ref{stratif} (iii) above, the first of these maps is a deformation retract and thus a homotopy equivalence.

We will now analyze the other inclusion map, which is $\Sigma_{\beta_0}\hookrightarrow  {\rm Gr}_k(\CC^n)$, see again
Thm.~\ref{stratif} (iii). We clearly have
\begin{equation}\label{zero}\Sigma_{\beta_0} = {\rm Gr}_k(\CC^n)\setminus \bigcup_{\beta \neq \beta_0} \Sigma_\beta.\end{equation}
Let $\beta^1_1, \ldots, \beta^1_{m_1}$ be the minimal elements of the indexing set $B$. By eq.~(\ref{osig}), the corresponding
$\Sigma$'s are closed in ${\rm Gr}_k(\CC^n)$. Prop.~\ref{prop1ton} and Lemma \ref{cs} imply that the map
$${\rm Gr}_k(\CC^n) \setminus \bigcup_{i=1}^{m_1}\Sigma_{\beta_i^1}\hookrightarrow {\rm Gr}_k(\CC^n) $$
is 3-connected. Let now $\beta^2_1, \ldots, \beta^2_{m_2}$ be the minimal elements of 
$B\setminus \{\beta^1_1, \ldots, \beta^1_{m_1}\}$. Again by eq.~(\ref{osig}), the spaces $\Sigma_{\beta_1^2}, \ldots, 
\Sigma_{\beta_{m_2}^2}$ are all closed in ${\rm Gr}_k(\CC^n) \setminus \bigcup_{i=1}^{m_1}\Sigma_{\beta_i^1}$. 
And by applying again Prop.~\ref{prop1ton} and Lemma \ref{cs}, the map
 $${\rm Gr}_k(\CC^n) \setminus \left(\bigcup_{i=1}^{m_1}\Sigma_{\beta_i^1}\cup \bigcup_{j=1}^{m_2}\Sigma_{\beta_j^2}\right)\hookrightarrow {\rm Gr}_k(\CC^n) \setminus \bigcup_{i=1}^{m_1}\Sigma_{\beta_i^1}$$
 is $3$-connected. The procedure can be continued and leads to the following chain of strict inclusions:
 $$B\supset B\setminus \{\beta_1^1, \ldots, \beta_{m_1}^1\}\supset B\setminus \{\beta_1^1, \ldots, \beta_{m_1}^1,
 \beta_1^2,\ldots, \beta_{m_2}^2\} \supset \ldots .$$
 The last set in the chain must be $\{\beta_0\}$, since $\beta_0$ is the greatest element of $B$ and hence not a minimal element of any subset of $B$ which contains it properly. 
We conclude that, in view of eq.~(\ref{zero}) above, the map
 $$\Sigma_{\beta_0} \hookrightarrow  {\rm Gr}_k(\CC^n) $$
 is a composition of several 3-connected maps and thus is itself 3-connected. 
\end{proof}

\begin{thebibliography}{99}
\bibitem[Ki]{Ki} F.~C.~Kirwan, {\it Cohomology of Quotients in Symplectic and Algebraic Geometry}, Mathematical Notes, vol.~31, Princeton Univ.~Press, New Jersey, 1984
\end{thebibliography}
\end{document}
